% --------------------------------------------------------------
% APA 7ma Edición - Plantilla para Tesis y Documentos Académicos
% --------------------------------------------------------------
\documentclass[12pt, a4paper, oneside]{article}

% Paquetes básicos
\usepackage[utf8]{inputenc}
\usepackage[T1]{fontenc}
\usepackage{graphicx}
\usepackage{amsmath}
\usepackage{amssymb}
\usepackage{booktabs}
\usepackage{setspace}
\usepackage{indentfirst}
\usepackage{url}
\usepackage[colorlinks=true, linkcolor=blue, citecolor=blue, urlcolor=blue]{hyperref}

% Configuración de página APA 7ma
\usepackage[right=1in, left=1in, top=1in, bottom=1in, includefoot]{geometry}

% --------------------------------------------------------------
% Portada APA 7ma
% --------------------------------------------------------------
\begin{document}
\begin{titlepage}
\centering
\vspace*{0.5in}
{\large\bfseries Título de la Investigación en Formato APA}\\[0.5in]
\vfill
\textbf{Autor 1}$^{1}$, \textbf{Autor 2}$^{2}$, \textbf{Autor 3}$^{3}$\\[0.3in]
$^{1}$Universidad, Departamento de Ingeniería\\
$^{2}$Universidad, Departamento de Ciencias\\
$^{3}$Universidad, Departamento de Investigación\\[0.5in]
{\small Director: Dr. Nombre del Director}\\[0.3in]
{\small Co-director: Dr. Nombre del Co-director}\\[0.5in]
{\small Ciudad, País}\\
{\small Año 2024}
\end{titlepage}

% --------------------------------------------------------------
% Página de aceptación
% --------------------------------------------------------------
\newpage
\thispagestyle{empty}
\vspace*{2in}
\begin{flushright}
\begin{minipage}{4in}
\small
\textbf{Fecha de sustentación:} dd de mes de aaaa\\[0.3in]
\textbf{Presidente del jurado:} Dr. Nombre\\[0.3in]
\textbf{Vocal 1:} Dr. Nombre\\[0.3in]
\textbf{Vocal 2:} Dr. Nombre\\[0.3in]
\end{minipage}
\end{flushright}

% --------------------------------------------------------------
% Abstract - Centrado, negrita
% --------------------------------------------------------------
\newpage
\section*{Abstract}
\addcontentsline{toc}{section}{Abstract}
{\centering\textbf{Abstract}\\[0.2in]}
\doublespacing
El resumen estructurado presenta de manera breve y concisa los aspectos
fundamentales de la investigación. Se estructura en: contexto, objetivo,
método, resultados y conclusión. El abstract debe contener entre 150 y 250
palabras.

\medskip
\textbf{Keywords:} palabra clave 1, palabra clave 2, palabra clave 3, palabra clave 4

% --------------------------------------------------------------
% Resumen en español
% --------------------------------------------------------------
\newpage
\section*{Resumen}
\addcontentsline{toc}{section}{Resumen}
{\centering\textbf{Resumen}\\[0.2in]}
El resumen en español presenta la misma información que el abstract,
estructurado de forma similar. Debe indicar el contexto de la investigación,
el objetivo principal, la metodología utilizada, los resultados obtenidos y
las conclusiones más relevantes. El rango permitido es de 150 a 250 palabras.

\medskip
\textbf{Palabras clave:} palabra clave 1, palabra clave 2, palabra clave 3

% --------------------------------------------------------------
% Tabla de contenidos
% --------------------------------------------------------------
\newpage
\tableofcontents

% --------------------------------------------------------------
% Lista de figuras y tablas
% --------------------------------------------------------------
\newpage
\listoffigures
\newpage
\listoftables

% --------------------------------------------------------------
% Cuerpo del documento
% --------------------------------------------------------------
\newpage
\selectfont
\doublespacing
\parindent=0.75in
\parskip=0in

% --------------------------------------------------------------
% Introducción - Sin número, negrita, centrada
% --------------------------------------------------------------
\section*{Introducción}
\addcontentsline{toc}{section}{Introducción}
La introducción presenta el problema de investigación, el contexto
teórico y los objetivos del estudio. En formato APA 7ma, la introducción
no lleva número y se identifica únicamente con el título centrado en negrita.

El problema de investigación se enmarca en el contexto de la disciplina
y se justifica su importancia (Apellido, Año). Los objetivos deben ser
claros, medibles y estar alineados con el problema planteado.

% --------------------------------------------------------------
% Marco Teórico o Antecedentes
% --------------------------------------------------------------
\section*{Marco Teórico}
\addcontentsline{toc}{section}{Marco Teórico}
El marco teórico presenta las bases conceptuales que sustentan la
investigación (Apellido, Año). Se estructura de forma lógica y
coherente, organizando los conceptos de lo general a lo específico.

\subsection*{Antecedentes de la Investigación}
Los antecedentes muestran investigaciones previas relacionadas con
el tema. Se presentan en orden cronológico inversal (Apellido, Año).

\subsection*{Bases Teóricas}
Las bases teóricas proporcionan el fundamento conceptual del estudio.

% --------------------------------------------------------------
% Metodología
% --------------------------------------------------------------
\section*{Metodología}
\addcontentsline{toc}{section}{Metodología}
La metodología describe el diseño de investigación, la población,
muestra, técnicas e instrumentos de recolección de datos, y el
procedimiento de análisis.

\subsection*{Diseño de Investigación}
Se especifica el tipo de investigación y su justificación.

\subsection*{Población y Muestra}
La población objetivo y la técnica de muestreo utilizada.

\subsection*{Técnicas e Instrumentos}
Los instrumentos de medición con su debida validación.

% --------------------------------------------------------------
% Resultados
% --------------------------------------------------------------
\section*{Resultados}
\addcontentsline{toc}{section}{Resultados}
Los resultados se presentan de forma objetiva, sin interpretación.
Pueden incluir tablas y figuras con caption centrado arriba.

\begin{table}[htbp]
\centering
\caption{Resultados del Análisis Descriptivo}
\begin{tabular}{lcc}
\toprule
Variable & Media & DE \\
\midrule
Variable 1 & 0.00 & 0.00 \\
Variable 2 & 0.00 & 0.00 \\
\bottomrule
\end{tabular}
\end{table}

% --------------------------------------------------------------
% Discusión
% --------------------------------------------------------------
\section*{Discusión}
\addcontentsline{toc}{section}{Discusión}
La discusión interpreta los resultados a la luz de la teoría
(Apellido, Año). Se comparan con estudios previos y se señalan
las limitaciones del estudio.

% --------------------------------------------------------------
% Conclusiones
% --------------------------------------------------------------
\section*{Conclusiones}
\addcontentsline{toc}{section}{Conclusiones}
Las conclusiones responden a los objetivos planteados. Se presentan
en orden correlativo y se sugieren investigaciones futuras.

% --------------------------------------------------------------
% Referencias - APA 7ma - Hanging indent
% --------------------------------------------------------------
\newpage
\singlespacing
\section*{Referencias}
\addcontentsline{toc}{section}{Referencias}
\parindent=-0.5in
\parskip=0in
\parskip=0.3em

% Referencias con hanging indent
\begin{thebibliography}{99}

% Artículo de revista con DOI
\vspace{0.3em}
Apellido, A. A., \& Apellido, B. B. (2024). Título del artículo.
\textit{Revista}, \textit{vol}(núm), páginas. https://doi.org/xxxxx

% Libro
\vspace{0.3em}
Apellido, A. A. (2023). \textit{Título del libro en cursiva} (2.$^a$ ed.).
Editorial.

% Capítulo de libro
\vspace{0.3em}
Apellido, A. A. (2023). Título del capítulo. En A. A. Apellido (Ed.),
\textit{Título del libro} (pp. xx-xx). Editorial.

% Tesis
\vspace{0.3em}
Apellido, A. A. (2023). \textit{Título de la tesis} [Tesis doctoral,
Universidad]. Nombre del repositorio. URL

% Página web
\vspace{0.3em}
Apellido, A. A. (2024). Título de la página. Nombre del sitio web.
https://ejemplo.com

\end{thebibliography}

% --------------------------------------------------------------
% Anexo A
% --------------------------------------------------------------
\newpage
\section*{Anexo A}
\addcontentsline{toc}{section}{Anexo A}
\textbf{Título del anexo}. Descripción del contenido del anexo.

% --------------------------------------------------------------
% Anexo B
% --------------------------------------------------------------
\newpage
\section*{Anexo B}
\addcontentsline{toc}{section}{Anexo B}
\textbf{Título del anexo}. Descripción del contenido del anexo.

\end{document}